% Frozen before the independent Aer rerun. Compile with xelatex.
\documentclass[11pt,a4paper]{article}
\usepackage{kotex}
\usepackage{amsmath,amssymb,booktabs,geometry}
\geometry{margin=25mm}
\title{I01 선검증 4 사전 프로토콜: 약한 합산 null의 독립 재시험}
\author{RELATIONAL ATLAS}
\date{2026년 10월 4일}
\begin{document}
\maketitle

\section{검정 수정 사유}
선검증 3의 층별 재배분 null은 각 표지에서 두 위상 블록의 성공확률이
같아야 한다. 하지만 균일 표지의 비제약 상태는 표지별 성공확률이
변해도 \emph{가중 합산} 확률만 정확히 같을 수 있다. 따라서 앞선
재배분 $p$값은 이 약한 합산 null에 대한 확증 판정에서 철회한다.
raw counts와 모든 사전 설정은 보존하고, 통계 방법을 고친 다음
독립 Aer seed를 새로 실행한다. 기존 raw counts의 재분석은
사후 진단으로만 표기한다.

\section{고정 실험과 약한 null}
물리·회로·표지별 shots 배분은 선검증 3과 동일하다:
$n=3,\beta=0.45,\beta_r=0.15$, $L=7$,
$E=E_0+0.37E_0^2/3$, $\tau=0,\pi/2$,
공통 위상 $0\to0.60$ rad, 균일/비균일 표지 및
비제약/$P=0$ 상태다. 각 축·상태·시계·위상 블록에
총 229376 shots, 합계 336회로다. CPU workers 6,
새 Aer seed 20261019를 쓴다. 여덟 사전 가족 검정의 null은
각 축 $i$에서
\begin{equation}
H_0:\quad \Delta_i=\sum_{t=0}^{6}a_t
\bigl[p_{1,it}-p_{0,it}\bigr]=0
\end{equation}
이다. 개별 표지 $p_{1,it}=p_{0,it}$은 요구하지 않는다.

\section{사전 통계량과 보정}
축별 표지 성공률 $\hat p_{b,it}$과 고정 shots 수 $N_t$에서
$a_t=N_t/229376$을 사용한다. 약한 null의 Wald 통계량은
\begin{align}
\hat\Delta_i&=\sum_t a_t(\hat p_{1,it}-\hat p_{0,it}),\\
\hat V_i&=\sum_t a_t^2\left[
\frac{\hat p_{0,it}(1-\hat p_{0,it})}{N_t}+
\frac{\hat p_{1,it}(1-\hat p_{1,it})}{N_t}\right],\\
T&=\sum_{i=1}^{3}\hat\Delta_i^2/\hat V_i.
\end{align}
주 판정의 null 분포는 관측 각 표지의 두 블록 확률 차이를 보존하면서
가중 합만 0으로 투영한
$p'_{0,it}=\hat p_{0,it}+\hat\Delta_i/2$,
$p'_{1,it}=\hat p_{1,it}-\hat\Delta_i/2$에서
독립 이항표본을 10000패널 생성해 얻는다. 투영 확률이 $[0,1]$
밖이면 멈추고 보고한다. 주 $p$값은 bootstrap 상측비율
$(1+\#\{T^*\ge T\})/10001$이고, $\chi^2_3$ 근사는
교차 점검용이다. 8검정 Bonferroni 문턱은 $0.00625$다.

정확 모형에서 균일/비제약 및 모든 $P=0$는 합산 null이고,
비균일/비제약은 두 $\tau$에서 양수다. 그러나 $\tau=\pi/2$의
정보는 이번 샷 예산에 비해 약해 미검출될 수 있다.
가족 문턱을 통과하지 않은 것을 영의 증거로 취급하지 않는다.
QPU는 쓰지 않고 발전연구 주제도 정하지 않는다.
\end{document}
